% Separate post-CP33 candidate. No change to the frozen B91/B92 proofs.
\subsection{The terminal determinant removes the last crossing}

Throughout this subsection the field is complete discretely valued
of mixed characteristic \((0,5)\), with algebraically closed residue
field, \(e_K=20\), and \(\mu _5\subset K\). Valuations on every
finite extension are normalized to have value group \(\mathbf Z\).
Lower breaks are \(b(g)=v(g\varpi-\varpi)-1\). The ramification
conductor of a nontrivial finite character means its last upper
break, without the additional one in the Artin-conductor convention.

\begin{lemma}[A lower bound for the actual cyclotomic conductor]
\label{lem:B-RD93-cyclotomic}
\status{Candidate; full proof and exact supplement pending independent audits.}
If \(\mu _{25}\not\subset K\), then
\(K(\mu _{125})/K\) is cyclic of degree 25 and its last upper
break is at least 20. More precisely, choose a Cohen coefficient
field \(K_0\subset K\), put \(K_1=K_0(\mu _5)\), and let
\(d=v_K(\mathfrak D_{K/K_1})\). Its last upper break is
\[
                              44-d\ge20.              \tag{B-RD93.1}
\]
After a tame extension of degree 24 the last upper break is
therefore at least 480.
\end{lemma}
\begin{proof}
The complete Cohen-ring argument used in B86 supplies \(K_0\)
with the given residue field. Since \(e(K_1/K_0)=4\),
the extension \(K/K_1\) is totally ramified of degree five.
For a uniformizer \(t\) of \(K\), its minimal polynomial is
Eisenstein of degree five. In its derivative the leading term
\(5t^4\) has valuation 24. Every other nonzero term has valuation
congruent to one of \(0,1,2,3\) modulo five, so cannot cancel
this term. Thus
\[
                                  d\le24.              \tag{B-RD93.2}
\]

Let \(T/K_1\) be the normal closure. Its group acts faithfully
and transitively on five letters. Because the residue field is
algebraically closed, the group is inertia; its wild subgroup is
normal and nontrivial. A five-subgroup of \(S_5\) has order at
most five, so the wild subgroup is the regular cyclic group of
order five. Its normalizer in \(S_5\) gives
\[
 \operatorname{Gal}(T/K_1)=C_5\rtimes C_m,\qquad m\mid4,
\]
where a nontrivial tame complement acts faithfully. There is just
one positive lower break, say \(b_T\), and \(T/K\) is tame of
degree \(m\). The different formula and transitivity give
\[
 5m-1+4b_T=(m-1)+md,
 \qquad u_T=\frac{b_T}{m}=\frac{d-4}{4}\le5.            \tag{B-RD93.3}
\]

Put \(C_2=K_0(\mu _{125})\) and \(C_1=K_0(\mu _{25})\).
The cyclic extension \(C_2/K_1\) has degree 25. With uniformizer
\(\zeta_{125}-1\), its order-25 elements have lower break four
and its nontrivial order-five elements have lower break 24:
the two differences have valuations five and 25, respectively.
Thus its upper breaks are four and eight.

The intersection \(T\cap C_2\) is \(K_1\). Indeed, a nontrivial
intersection would contain \(C_1\), and give a cyclic quotient
of order five of \(\operatorname{Gal}(T/K_1)\). This is impossible
when \(m>1\), since the faithful tame action kills its wild
subgroup in the abelianization. When \(m=1\), one has \(T=K\),
and the intersection would force \(K=C_1\), contrary to the
hypothesis. In particular \([KC_2:K]=25\).

The last upper break eight of \(C_2/K_1\) is strictly above
the last upper break of \(T/K_1\). Upper-quotient compatibility
and the subgroup Herbrand conversion therefore give the last
upper break of \(TC_2/T\) as
\[
 \psi_{T/K_1}(8)
 =m u_T+5m(8-u_T)=m(44-d).                            \tag{B-RD93.4}
\]
This is an equality: the order-five last cyclotomic subgroup
is still present after the other factor has become trivial.
The extension \(T/K\) is tame of degree \(m\), and is disjoint
from the cyclic 25-extension \(KC_2/K\). Tame base change
multiplies upper breaks by \(m\), proving (B-RD93.1).
The final assertion follows by the same tame conversion.
All subgroup and quotient conversions here are those of the
complete Herbrand proof already used in B88--B92; no finite
residue-field hypothesis is introduced.
\end{proof}

\begin{proposition}[The actual second division field and a five-dimensional quotient]
\label{prop:B-RD93-second-field}
\status{Candidate; full proof and exact supplement pending independent audits.}
Suppose the surviving B92 crossing \((B_1,B_2)=(5,45)\) occurs.
Let \(F/L\) be its common actual 125-division field, \(F_2/L\)
its original 25-division field, and \(F_Q/L\) the common
25-division field of the actual high quotients. Put
\(E=F_2F_Q\). Then
\[
 \begin{array}{c|c|c}
 \text{extension}&\text{lower breaks}&\text{upper breaks}\\\hline
 F_2/L&24,234,384,44634&24,66,72,426\\
 E/L&6,96,1146,1896,223146&6,24,66,72,426.
 \end{array}                                         \tag{B-RD93.5}
\]
The degrees are respectively \(5^4\) and \(5^5\).

Write the full parahoric matrix of B92 as
\[
 \rho=\begin{pmatrix}1+5a&b\\25c&1+5d\end{pmatrix},
 \quad a,d\in\mathbf Z/25,\quad b\in\mathbf Z/125,
 \quad c\in\mathbf F_5.
                                                               \tag{B-RD93.6}
\]
Set \(a_0=a\bmod5\), \(d_0=d\bmod5\), and \(u=b\bmod5\).
Let \(R,A,D,C\) be the cyclic fields of \(u,a_0,d_0,c\),
respectively. They have breaks \(66,72,24,6\) over \(L\).
For \(S=RA\), and for an actual primitive 25-point
\(P_\ell\) with \(5P_\ell=e_2\), put \(E_\ell=L(P_\ell)\).
Then
\[
 \begin{array}{c|c|c|c}
 \text{extension}&\text{degree}&\text{lower breaks}&\text{upper breaks}\\\hline
 E/S&125&6,96,223146&6,24,8946\\
 E/E_\ell&25&6,1896&6,384.
 \end{array}                                         \tag{B-RD93.7}
\]
Finally
\[
 U=\operatorname{Gal}(F/E)
 =\left\{I+25\begin{pmatrix}a_1&b_2\\0&d_1\end{pmatrix}\right\}
 \simeq\mathbf F_5^3.                                \tag{B-RD93.8}
\]
\end{proposition}
\begin{proof}
B92 proves equality with the entire group (B-RD93.6), not only
containment, and identifies all four cyclic fields. The kernel
of the original mod-25 image is all of \(M_2(\mathbf F_5)\)
at the terminal level. Over \(R\), the original mod-25 image is
the additive upper-triangular three-space, with coordinates
\(a_0,d_0,b_1\), where \(b=5b_1\pmod {25}\).

The actual point field \(E_\ell/R\) has degree 25 and is
Galois, with translation coordinates \(b_1,d_0\). By the
common-high calculation of B92 for \((b_1,b_2)=(1,9)\), its
lower breaks are 24 and 9234, hence its upper breaks are
24 and 1866. Its last line is the upper-entry line, so the
\(d_0\)-character has break 24 and the \(b_1\)-character has
break 1866. The actual \(a_0\)-character over \(R\) has break
\(66+5(72-66)=96\), by the two distinct cyclic fields
\(A/L\) and \(R/L\). The three distinct conductors give the
upper breaks \(24,96,1866\) of \(F_2/R\), and thus its
lower breaks \(24,384,44634\).

Here is the complete passage back to \(L\), without assigning
breaks to arbitrary lifts. The nonidentity elements of
\(H=\operatorname{Gal}(F_2/R)\) have break 24, 384, 44634
with multiplicities 100, 20, 4. In any nontrivial coset of
\(H\), let \(M\) be its maximal lower break. Choosing a lift
with this break shows that the other breaks are the minimum
of \(M\) and the corresponding subgroup break. When the
two are equal, a larger value is impossible by maximality,
so the same formula still applies. The exact quotient average
therefore says
\[
 \frac{100\min(24,M)+20\min(384,M)+4\min(44634,M)+M}{125}
 =66.                                                 \tag{B-RD93.9}
\]
The left side is strictly increasing; its unique solution is
\(M=234\). Each of the four cosets consequently has 100
elements of break 24 and 25 of break 234. This proves the
first row of (B-RD93.5).

The character \(c\) is jointly independent of the entire
original mod-25 image: the terminal lower matrix entry varies
freely while that image is fixed. Thus
\(E=F_2C\), and its group is the direct product of the
original mod-25 group with the cyclic group of order five.
The latter character has upper break six, below all the
breaks of \(F_2/L\). Combining these actual quotient
filtrations gives the second row of (B-RD93.5).

For completeness, both relative fields in (B-RD93.7) can be
read directly in this group. The stabilizer of \(S=RA\)
is defined by \(u=a_0=0\); it has free coordinates
\(d_0,b_1,c\). The corresponding breaks are 96, 223146,
and six, giving that row. The stabilizer of \(E_\ell\)
is defined by \(u=b_1=d_0=0\); it has free coordinates
\(a_0,c\). Its lower breaks are 1896 and six. Subgroup
lower numbering and then the upper conversion give precisely
(B-RD93.7). These are actual stabilizers of the stated fields.
Intersecting the kernels of the original mod-25 image and
the quotient mod-25 image gives exactly (B-RD93.8).
\end{proof}

\begin{lemma}[Actual primitive points bound both terminal diagonal characters]
\label{lem:B-RD93-terminal-bound}
\status{Candidate; full proof and exact supplement pending independent audits.}
In the situation of the preceding proposition, the characters
\(a_1,d_1\) of \(U\) satisfy
\[
 \operatorname{cond}_E(a_1)\le275496,
 \qquad \operatorname{cond}_E(d_1)\le70296.             \tag{B-RD93.10}
\]
These bounds use the actual 125-division fields and do not
require a 625-torsion comparison.
\end{lemma}
\begin{proof}
Choose the integral tame eigenbasis used in B92. A high inverse
of the first-level point \(e_1\) in the first model lies in
the unique high inverse coset, of size five, with valuation 19.
This coset is stable under the tame generator after the
appropriate eigenvalue rescaling. Since the other residual
tame eigenvalue is distinct, its prime-to-five eigenprojector
selects a lift in the \(e_1\) eigenline without leaving that
coset. Thus we may take the primitive 25-point in this
eigenline to be high in the first model. Its valuation in
the second model is three, since all 25 inverses there have
that valuation. Choose an actual primitive 125-point \(P_h^+\)
lifting it in the same tame eigenline.

The inverse Newton polygons for \(P_h^+\) are single segments
in both models, with valuations \(19/25\) and \(3/25\)
over \(L\). Indeed their target valuations are 19 and three,
while the degree-five coefficient valuations are 20 and 180;
both degree-five vertices lie strictly above the corresponding
degree-25 segment.

Let \(E_h^+=L(P_h^+)\) and \(T_h=E_h^+R\). The actual
matrix stabilizer of \(P_h^+\) is \(a=c=0\), so
\([E_h^+:L]=125\); the character \(u\) still varies freely
on that stabilizer, giving \([T_h:L]=625\). The field
\(S=RA\) is contained in \(T_h\), and
\[
 \operatorname{Gal}(T_h/S)=\mathbf F_5^2
 \quad\text{with translation coordinates }(a_1,c).
                                                               \tag{B-RD93.11}
\]
For example this follows directly by restricting (B-RD93.6)
to \(u=a_0=0\); its action on the first column is addition
by \(25(a_1e_1+ce_2)\), and the kernel is normal. The
quotient killing the high translation line \(c=0\) is
the actual field \(CS/S\), with break six: both breaks
66 and 72 of \(S/L\) are larger than six.

In \(T_h\), the second-model coordinate has valuation 75.
Every nonzero element of the high translation line has
displacement valuation \(625\cdot75=46875\). The fixed-Cohen-
representative inequality, valid even when the coordinate
valuation is divisible by five, gives
\[
                    b(h)\le46875-75=46800.             \tag{B-RD93.12}
\]
If the high line is the last ramification line, the first
lower break is six, and (B-RD93.12) bounds the last upper
break by \(6+(46800-6)/5=46824/5\). If it is not last,
the character \(c\) is nonzero on the last line, so its
break six is already the largest upper break. A one-break
extension has the same conclusion. In every case the
\(a_1\)-character of (B-RD93.11) has conductor at most
\(46824/5\).

On restriction to \(E\), this is precisely the character
\(a_1\) in (B-RD93.8). The subgroup/quotient Herbrand formula
bounds its conductor by \(\psi_{E/S}(46824/5)\).
Using the first row of (B-RD93.7) gives
\[
 223146+125\left(\frac{46824}{5}-8946\right)=275496.
                                                               \tag{B-RD93.13}
\]
This is an upper bound and does not assume disjointness of
the two high-point fields; in fact \(T_h\cap E=SC\).

For the other diagonal character, choose an actual primitive
125-point \(P_\ell^+\) with \(5P_\ell^+=P_\ell\).
The full matrix group gives
\[
 [L(P_\ell^+):L]=3125,\qquad
 \operatorname{Gal}(L(P_\ell^+)/E_\ell)=\mathbf F_5^2,
                                                               \tag{B-RD93.14}
\]
with translation coordinates \((b_2,d_1)\). The stabilizer
at level 125 is \(b=d=0\); at level 25 it is
\(b\equiv0\pmod {25}\), \(d\equiv0\pmod5\).
The former is normal in the latter, and the displayed
quotient is their actual action on the second column.

In the first model the point \(P_\ell\) has valuation
\(1/25\) over \(L\), and all its inverses have valuation
\(1/625\). Thus in \(L(P_\ell^+)\) the coordinate value
is five. Every translation with \(d_1\ne0\) is outside
the canonical high line and has displacement value 3125.
The same unconditional inequality gives \(b(g)\le3120\)
for every such element. The exact quotient average over
the high line therefore bounds the cyclic \(d_1\)-quotient
break by 3120. Restriction to \(E\), using the second row
of (B-RD93.7), bounds its conductor by
\[
 \psi_{E/E_\ell}(3120)
 =1896+25(3120-384)=70296.                              \tag{B-RD93.15}
\]
Here \(E\cap L(P_\ell^+)=E_\ell\), as is also immediate
from the two independent terminal column coordinates.
All point fields, kernels and restrictions in these two
calculations are attached to the original actual 125-module.
\end{proof}

\begin{theorem}[The last residual-wild crossing is impossible]
\label{thm:B-RD93-crossing}
\status{Candidate; full proof and exact supplement pending independent audits.}
The remaining crossing \((B_1,B_2)=(5,45)\) of B92 cannot
occur. Consequently every remaining residual-wild actual
125-comparison at \(p=5,e=20\) with different finite special
characters has both original parameters above 20 and the
two layer shapes \((3,4)\) or \((4,4)\) of B91.
\end{theorem}
\begin{proof}
The finite determinant is trivial on the common good field
\(L\), so the determinant of the actual 125-module on
\(\operatorname{Gal}(F/L)\) is the cyclotomic character
(or its inverse, with identical conductor). Its mod-25
character is \(a_0+d_0\). The two summands have breaks
72 and 24, respectively, so this character is nontrivial
and has break 72. Thus \(\mu _{25}\not\subset K\).

By Lemma~\ref{lem:B-RD93-cyclotomic}, the actual cyclic
extension \(L(\mu _{125})/L\) has last upper break
\(v\ge480\). This is strictly above the largest upper
break 426 of \(E/L\). Consequently the last cyclotomic
subgroup persists in the subgroup over \(E\), and the
conductor of the restricted determinant is exactly
\(\psi_{E/L}(v)\). In particular it is at least
\[
 \psi_{E/L}(480)=223146+3125(480-426)=391896.             \tag{B-RD93.16}
\]

On the actual kernel (B-RD93.8), however,
\[
                  \det\rho=1+25(a_1+d_1)\pmod {125}.
\]
The conductor of a sum of two finite characters is at most
the larger of their conductors. Lemma~\ref{lem:B-RD93-terminal-bound}
therefore bounds this same restricted determinant by
\(275496<391896\), a contradiction. This comparison is
between the same actual character of \(\operatorname{Gal}(F/E)\),
not an identification of abstract cyclic quotients.
The final statement now follows from B91 and B92.
\end{proof}

\begin{remark}[Scope and retained checks]
\label{rem:B-RD93-scope}
\status{Candidate scope; no complete first-wild least-index assertion.}
This proposed exclusion addresses the original residual-wild
crossing only. The both-large alternatives of B91 and the
irreducible-residual sectors are not excluded or realized.
The unrestricted first-wild least index consequently remains
in \(\{2,3\}\) unless those remaining actual comparisons are
separately settled. No 625-level comparison is assumed after
passing to any quotient. The normal-closure/different argument
is independent of local class field theory and of finite
residue fields.

The complete invoked B60 fixed-representative argument,
B86 Cohen-field construction, B88--B90 actual field and
Herbrand conversions, and B91--B92 quotient/stabilizer proofs
are the dependencies. The explicit degree-five normal-closure
and cyclotomic argument is supplied above. No fresh external
original-source opening is claimed here. The retained original
source conventions are the full Cohen proof and the different,
subgroup and quotient Herbrand proofs already checked in B86--B90
\cite{Stacks032A,ColmezCorpsLocaux}.

The standalone exact supplement
\path{code/reduction_descent/verify_residual_wild_last_crossing.py}
checks all 35 displayed normal-closure arithmetic cases without
asserting their existence, enumerates all 390,625 actual matrices
of the full parahoric group and reconstructs their ten point/field
images, stabilizers and the two intersections. It independently
performs the rational Herbrand conversions, the coset average and
all 125 terminal determinant checks. The bounds 275,496 and
70,296 are strictly below the cyclotomic lower bound 391,896,
with gap 116,400. The first author run passed at
03:48:13.668022 UTC on 20 November 2026 in
0.6302177960023982 seconds. No executed computation failed.
The full source, output, pre-certificate version and exploratory
provenance note are retained. These finite checks supplement the
actual field proof; no finite arithmetic row is used as an
existence assertion.
\end{remark}
